\chapter*{译校说明}
\markboth{译校说明}{译校说明}
\addcontentsline{toc}{chapter}{译校说明}

\noindent\small
\textbf{底本。}本书译自 Paul Seidel 的讲义 \emph{Fukaya Categories and Picard--Lefschetz Theory}（欧洲数学会，2008；作者公开的 PDF 版本，共 334 页）。翻译底本的原书页码与 PDF 页码相差 8：PDF 第 9 页对应原书正文第 1 页。书中脚注与译者注所称“原书 p.~$N$”均指原书\textbf{印刷页码}，与 PDF 阅读器的页码不同。

\textbf{初版与本次校订。}初版（修订版 01）据英文原文译出，止于第 20 节 20.7 末，共 243 页，未加译者注。本次校订把全书 82 个章节文件按原书页序分成 12 段，逐段与英文原文对照通读，公式、交换图与符号约定另据原书页面图像逐式核对；遇原书排印可疑处，先核对页面图像，再决定改译文还是存疑。校订只改真正的误译、漏译、术语不统一与公式/指标错误，不作风格性改写；原书自身的笔误保留原译的正确写法，并在译者注与《原书笔误登记》中说明，不回改。

\textbf{译者注。}全书译者注统一排在书末“译者注”附录，正文只在相应句末加数字上标，不改动正文语义。每条注以“【原书 p.~$N$】”开头，先给页码，再给注文，注末给依据。依据分三级：A 级为原书前后文、作者他文、经典教材与原始论文；B 级为 nLab、Stacks Project、EMS 等可靠百科或数据库，均经实际打开核对并注明访问日期；此类注只作概念与背景说明，凡涉及具体定理归属或数值者，另举原书前后文、经典教材或原始论文为据（A 级），不单凭百科立论。无外部依据、仅由上下文推得者，明确写作“译者按（上下文推断，无外部依据）”。凡未能核实者不入注，宁缺勿滥。全书共 187 条译者注，A 级 125 条、B 级 51 条、C 级 11 条。

\textbf{术语。}全书术语以书末《术语表》为准。较初版有两处统一：Lefschetz thimble 由“指套”改为“顶针”（thimble 为顶针，finger cot 才是指套）；operad 由“运算胚”改为“操作子”。前者是初版沿用的日文译法，后者“胚”字在中文数学中另有 germ 的含义，均易致误，故改。

\textbf{使用建议。}初次阅读可先读“引言”与“读书报告”，把握定理 A、B 与推论 C 三条主线；再按章节顺序推进。译者注按需查阅，不必先读。书中少数记号（如用 \S{} 形符号表示 Koszul 符号）与原书一致，遇符号歧义时以注文说明为准。

\textbf{免责。}译者注是理解辅助，不是原文的一部分；凡注文与原文冲突，以原文为准。译文与注文中的疏漏由译者负责，欢迎指正。
\normalsize
